\cotefile{82}
\scanpage{50}
% @ lines 1–6 / text
Mais il y a dans $Q$ un \unsure{autom.} involutif canonique, qui correspond à la symétrie des deux facteurs $X\times X$ dans (6). Il est induit par un antiautomorphisme \inserted{can.} de l'Algèbre $M$ obtenu \inserted{d'abord $\mathrm{loc}^{\mathrm t}$} en choisissant \deleted{\unsure{loc$^{\mathrm t}$ une base}} une \deleted{\unsure{$\omega$}} base $\omega$ de $\det\check V$\inserted{$=\bigwedge^2\check V$}, d'où un isom \deleted{\illegible{1}}
% @ line 7 / display
\[
 V\xrightarrow{\deleted{\sim}\,\alpha}\check V
 \qquad\text{composé de}\quad V\longrightarrow\check V\otimes\det V\xrightarrow{\omega}\check V
 \tag*{(13)}
\]
% @ line 8 / text
ce qui permet de définir
% @ line 9 / display
\[
 u\longmapsto u'=\alpha^{-1}({}^{t}u)\alpha\tag*{(14)}
\]

% @ lines 10–12 / text
Mais changer de $\omega$ revient à multiplier $\alpha$ par une \deleted{section} unité $\lambda\inG_m(S)$, donc \deleted{\unsure{$u$ ch.}}\inserted{$\alpha^{-1}$ devient $\lambda^{-1}\alpha$} et $u'$ ne change pas. Par construction
% @ line 13 / display
\[
 (uv)'=v'u' \tag*{(15)}
\]

% @ line 14 / text
En fait, le calcul donne
% @ line 15 / display
\[
 u'=\operatorname{Tr}u.1-u\tag*{(16)}
\]
% @ line 16 / text
plus richement on a
% @ lines 17–18 / display
\[
 \left\{\begin{aligned}u+u'&=(\operatorname{Tr}u)\,1\\ uu'&=(\det u)\,1\end{aligned}\right. \quad,\tag*{(17)}
\]
% @ lines 19–21 / text
et il est immédiat que $u\longmapsto u'$ échange $\operatorname{Ker}$ et $\operatorname{Im}$ pour un $u$ tel que $P(u)=0$, \unsure{ce|$u$} non nul sur les fibres....

% @ lines 22–25 / text
Quand on voit ce que devient cette involution de $M$ et $Q$ en tenant compte de (12), on voit qu'à tt couple $(A,\rho)$ elle associe $A$ ($=A'$ !) et $\rho'$ défini par
% @ line 26 / display
\[
 \rho'(u)=\rho(u')\tag*{(18)}
\]
% @ line 26, lower right / margin note
\margin{(\uline{NB} $u'|A$ est un \unsure{autom.|automorphisme}, car $A$ comm.)}
% archive sheet number: 49
\scanpage{51}
% @ line 1 / text
Considérons l'hom
% @ lines 2–3 / display
\[
 Q\longrightarrow\check{\mathbb P}(\unsure{\mathfrak g|g})
 =\underset{\substack{\shortparallel\\\Sigma}}{\mathbb P(\widetilde{\mathfrak g})}
 \quad(=\ \text{Schéma des ss-Algèbres spéciales de }M
 \tag*{(19)}
\]
% @ lines 4–6 / text
C'est évidemment un \uline{épi}, et si $A_\Sigma$ est ``la sous-algèbre spéc. universelle'' de $M$ -- ou \unsure{plutôt}, de $M_\Sigma$ -- on trouve donc un iso. can.
% @ line 7 / display
\[
 Q\simeq\operatorname{Spec}(A_\Sigma)\tag*{(20)}
\]
% @ left below (20) / text
\noindent\shortstack{\rule{1.7cm}{0.4pt}\\[-0.55em]\rule{1.7cm}{0.4pt}}

% @ lines 8–15 / text
Ceci étant, $A$ étant fixée, et $\underline L,\underline M$ étant des sous-fibrés inv. de $\widetilde{\mathfrak g}$ correspondants à des sections de $X$ (i.e. sur lesquels $q$ s'annule) à quelle condition le $\underline J$ \deleted{\unsure{que}}\inserted{$=L^\perp\otimes M$} $\subset M$ section de $Q$ qui leur correspond est-elle $\subset A$ ? (Les couples $\underline L,\underline M$ en question correspondent donc exactement aux sections de $Y=\operatorname{Spec} A$ sur $S$.)\unsure{.}

% @ line 16, first cancelled line / deletion
\deleted{Je dis qu'il \unsure{faut s.} que $L^{\unsure{\epsilon|{}}}\subset A^\perp$}
% @ lines 17–19, crossed proposition / text
\deleted{\uline{Prop} Conditions équivalentes, pour $A,L,M$ donnés} \deleted{\unsure{(i.e.}} \deleted{\unsure{$\Leftrightarrow L\subset$}}
% @ lines 18–19, bracketed hypotheses / display
\[
 \deleted{\unsure{\Leftrightarrow}\left\{\begin{aligned}
 &A\ \text{ss-Alg. spéciale de }M\\
 &L,M\ \text{ss-fibrés }\inserted{\text{inv de rang }1\text{ de }V}\ 
 \deleted{\text{satisfaisant }\operatorname{Tr},\det|L,M=0}
 \end{aligned}\right.}
\]
% @ lines 20–21, paired conditions / display
\[
 \deleted{\left.\begin{aligned}
 a)\quad &\underline L\subset A^\perp\quad(\text{i.e }\operatorname{Tr}\underline LJ=0\quad\text{si } A=\mathcal O_S\oplus J)\\
 b)\quad &\underline M\subset A^\perp\quad(\text{i.e. }\operatorname{Tr}\underline MJ=0\quad\text{————})
 \end{aligned}\right\}}
\]
% @ line 22 / display
\[
 \deleted{c)\quad L^\perp\otimes M\subset A}
\]
% @ lines 23–25, end at lower edge / text
\deleted{De plus, pour tt $L$\inserted{$\in X(S)$} \deleted{satisfaisant \unsure{a)}}, \inserted{$L\subset A^\perp$,} il existe un et un seul $M\in X(S)$ satisfaisant c) [ donc}
% archive sheet number: 50
\scanpage{52}
% @ lines 1–4 / text
\uline{Proposition} Soient $A$ une ss-Alg. spéciale, et $L$ un sous-fibré inv. de $V$, $\underline L=L^\perp\otimes L\subset M$ le sous-fibré isotrope de $\widetilde{\unsure{\mathfrak g|g}}$ associé. Conditions équivalentes
% @ line 5 / display
\[
 a)\quad \underline L\subset A^\perp
\]
% @ lines 6–7 / text
b) $A$ invarie la droite $L$ (\uline{NB} si $A=\mathcal O_S+J$, il revient au $=$ bien sûr que $J.L\subset L$)
% @ line 7, smaller line below b / text
c) $\exists$ ss-fibré inv. $J$ de $A$ tel que $J.L=0$
% @ line 8 / text
d) $\exists M\subset V$ ss fibré vectoriel \unsure{de rg|droite} tel que
% @ line 9 / display
\[
 \deleted{\unsure{M}}\qquad L^\perp\otimes M\subset A
\]

% @ lines 10–13 / text
De plus, sous ces conditions, le $M$ du c) est unique, et l'unique hom $A\longrightarrow\mathcal O_S$ de noyau $L^\perp\otimes M$ \deleted{\unsure{provient de}} est aussi l'hom $A\xrightarrow{u}\mathcal O_S\simeq\underline{\operatorname{End}}(L)$ déduit de b).
% @ lines 14–17 / text
Enfin, si $J$ est comme dans c), ce noyau est aussi égal à $J$. Et pour $A$ fixé, la corr. $L\mapsto u$ est une bijection entre l'ens des $L$ satisfaisant à a), et l'ens des sections \unsure{d'alg} $A\to\mathcal O_S$

% @ lines 17–19, proof beginning in interline / text
\uline{Dém} La question étant locale, \unsure{sps} $L$ \deleted{\unsure{engendré}} et $L^\perp$ munis de bases $x,x'$, \inserted{\unsure{de} $x'\otimes x$ base de $\underline L$.} si $u$ est \deleted{un} un endom de $V$ i.e section de $M$, on a
% @ line 20 / display
\[
 \operatorname{Tr}u\,x'\otimes x=\langle ux,x'\rangle
\]
% @ lines 21–25 / text
qui est nul ssi $ux$ est orth. à $x'$ i.e. $u(L)\subset L$, d'où a)$\Longleftrightarrow$b). Evidemt si $u:A\to\mathcal O_S\simeq\underline{\operatorname{End}}(\underline L)$ est l'hom correspondant, il est épi donc son noyau $\underline J$ est un ss-fibré inv. tel que $J.\underline L=0$, donc b)$\Rightarrow$c).
% @ lines 26–28 / text
Inv. pour un tel $\underline J$, on sait \inserted{(\unsure{v. leçon.|v. leçons})} que $\mathcal O_S\xrightarrow{\sim}\underline{A/J}$ (car $A/J$ $\mathcal O_S$-Alg. libre de rg $1$) donc $A\simeq J+\mathcal O_S$ et $L$ est stable sous $A$, d'où b)$\Longleftrightarrow$c)
% archive sheet number: 51
\scanpage{53}
% @ lines 1–2 / text
De plus, cet argument montre aussi que $J$ est uniquement déterminé comme $\operatorname{Ker}(A\to\operatorname{End}(\underline L))$.
% @ lines 3–8 / text
De plus, sous les conditions équivalentes a) b) c) \deleted{on aura évidemment $\det|\underline J=0$ donc $J$ \unsure{se}.} soit loc. $u$ une base de $\underline J$, \deleted{\unsure{donc est soit|donc est sur} de la forme $\underline L\otimes$} de sorte que $L=\operatorname{Ker}u$, \unsure{soit|donc} $M=\operatorname{Im}u$, c'est un ss-fibré vectoriel inv. tel que \deleted{$L^\perp\otimes M$} $L^\vee\otimes M=\mathcal O_Su=\underline J\subset A$, d'où d).
% @ lines 8–13 / text
D'ailleurs, si $M$ est un ss-fibré inv. de $V$ tel que $L^\vee\otimes M\subset A$, évidemt $L^\vee\otimes M$ est un ss-fibré inv. de $A$ nul sur $L$, donc c'est $J$, ce qui montre en $=$ temps que $M$ est uniquement déterminé comme $\operatorname{Im}u$ ($u$ base loc$^{\mathrm{le}}$ de $J$).
% @ lines 14–17 / text
Reste à prouver d)$\longrightarrow$b), or \unsure{d)} implique $A=\mathcal O_S+J$ (où on pose $J=L^\perp\otimes M$) \deleted{et} $J$ étant nul sur \deleted{\unsure{$\mathcal O$}} $\underline L$ à fortiori l'invarie cqfd.

% @ line 18, overwritten heading / deletion
\deleted{\uline{Corollaire}}

% @ lines 19–23 / text
Se donner un $A$ (ss fibré vectoriel de rg $2$ de $M$ contenant $1$\inserted{i.e.} $\mathcal O_S.1$ -- un tel $A$ est néc une ss-alg comm !) ou se donner $A^\perp=H\subset\widetilde{\unsure{\mathfrak g|g}}$, un sous fibré vectoriel de rang $2$ \inserted{i.e. corg $1$,} de $\widetilde{\mathfrak g}$, revient au $=$
% @ line 24, lower edge / display
\[
 \operatorname{AlgSpec}(M)\simeq
 \underbrace{\mathbb P(\widetilde{\mathfrak g})\;(\simeq\check{\mathbb P}(\mathfrak g))}_{\check P}
 \simeq\left\{\begin{array}{l}
 \inserted{\text{schéma des}}\\[-2pt]
 \text{fibrés en droites)}\\
 \text{dans }P\mathbin{:}\check{\mathbb P}(\widetilde{\mathfrak g})=\mathbb P(\mathfrak g)
 \end{array}\right.
 \tag*{(21)}
\]
% archive sheet number: 52
\scanpage{54}
% @ line 1 / text
L'\deleted{application} morphisme)
% @ lines 2–3 / display
\[
 \begin{aligned}
 D&\longmapsto D\cap X=Y\\
 \check P=\operatorname{Dr}(P)&\xrightarrow{\sim}\operatorname{Div}_2^+(X)
 \end{aligned}\tag*{(22)}
\]
% @ lines 4–6 / text
est un iso, qui associe à \deleted{\unsure{tt}} $H\subset\widetilde{\unsure{\mathfrak g|g}}$ le diviseur $Y\subset X$ des \deleted{droites}\inserted{ss-fibrés inv. $\underline L$} dans $\widetilde{\mathfrak g}$ qui \uline{à la fois} sont dans $H$, et satisfont $q|\underline L=0$.
% @ lines 7–10 / text
\deleted{C'est donc} Comme schéma sur $S$, il représente donc le sous-foncteur de $\unsure{\mathbb P^{*}|\check{\mathbb P}}(H)=\check{\mathbb P}(H)$ défini par l'annulation de $q$. On trouve donc, grâce à la \deleted{\unsure{proposition}}\inserted{prop. précédente}, un iso. can.
% @ line 11 / display
\[
 Y\simeq\operatorname{Spec} A\tag*{(23)}
\]

% @ lines 12–16 / text
Dans l'interprétation $X\xleftarrow{\sim}\mathbb P(V)$, via $L\rightarrow\underline L$, on voit que $Y$ est le sous-schéma de $\mathbb P(V)$ formé des $L$ tels que $AL\subset L$ ou encore (si $A=\mathcal O_S1+\mathcal O_Su$) tel que $uL\subset L$, ou encore, si $u$ est inv. i.e. section de $\check{\mathbb V}(M^*)$
% @ below line 16, left / deletion
\deleted{\illegible{1}}
% @ lines 17–18 / display
\[
 Y=X^u\qquad\begin{aligned}
 &\text{sous-schéma des}\\[-2pt]
 &\text{invariants sous }u
 \end{aligned}\tag*{(24)}
\]
% @ lines 19–24, continues on next page / text
On retrouve réciproquement $A$ quand on connaît $Y$, en notant que \deleted{l'opération de}\inserted{l'isomorphisme} (23) est fonctoriel, donc compatible aux autom provenant de $u\in\deleted{\unsure{A}}\Gamma(M^*)$ (i.e. ici $\{u:V\to V)$, en faisant opérer $u$ par autom int. sur $M^*$, et de la façon évidente sur $X$ \unsure{$\leftarrow$|\ldots}
% archive sheet number: 53
\scanpage{55}
% @ lines 1–2 / text
On trouve, en notant que tt ne dépend que de la section correspondante de $G=\check{\mathbb V}(M^*)/G_m$ :
% @ lines 2–3 / text
\uline{Prop} Pour que $g=\bar u\in G(S)$ invarie $Y$, il f. et s. qu'il invarie $A$ i.e. que \deleted{\unsure{$\operatorname{int}(u)A=A$}.}
% @ line 4 / display
\[
 \operatorname{int}(u)A\ \bigl(\overset{\mathrm{dfn}}{=}uAu^{-1}\bigr)=A
\]
% @ lines 5–8 / text
Quand il en est ainsi, l'opération de $g$ sur $Y$ \deleted{est} s'identifie via (23) à l'op. de $g$ \inserted{(ou de $u$)} sur $\operatorname{Spec}(A)$ -- donc elle est triviale ssi $g$ opère trivialement sur $A$ i.e ssi $u\in\Gamma(S,A)$.

% @ lines 9–11 / text
Bien entendu, ceci caractérise $A$ \deleted{\unsure{via|comme}} (engendré comme sous-fibré vectoriel par les sections locales de $A^*$ ...).

% @ line 12 / text
À \deleted{\unsure{tout}} $A$, associons le ss-grpe
% @ lines 13–14 / display
\[
 \underset{\substack{\shortparallel\\\check{\mathbb V}(A)^*/G_m}}{\unsure{\mathscr J|\mathscr T|\mathcal T}}\subset G\tag*{(25)}
\]
% @ lines 15–20 / text
c'est un sous-grpe \uline{commutatif} \uline{lisse} \uline{connexe} de dim \uline{relative} $1$. Bien entendu, la correspondance $A\mapsto\mathscr J$ est injective (\deleted{\unsure{$\mathbb V$}} $A$ étant engendré par \deleted{\unsure{l'}} les sections de $\check{\mathbb V}(A)^*$, qui est l'image inverse de $\mathscr J$ par l'hom can $\check{\mathbb V}(M)^*\longrightarrow G$.
% @ lines 21–24, continues / text
D'autre part, \inserted{si $2$ est inv. sur $S$,} tout $\mathscr J$ ayant les propriétés énumérées (sans $=$ la \unsure{commutati...|commutativité...|connexité...}) \deleted{-- appelons les les ss-grpes \uline{réguliers} de $G$ -- proviennent de tels $A$} -- i.e. de ss. fibrés inv
% archive sheet number: 54
\scanpage{56}
% @ lines 1–3, top-edge return to earlier sentence / text
de $\unsure{\mathfrak g|g}=\operatorname{Lie}G$ $+$ \inserted{sont de la forme $\check{\mathbb V}(A^*)/G_m$ -- i.e. sont des ss-grpes ``\uline{jacobiens}''.} En effet, considérons $\mathfrak h=\operatorname{Lie}\unsure{\mathscr J|\mathscr T|\mathcal T}$, qui est un tel ss-fibré, ($\mathscr J$ lisse de dim rel. $1$)
% @ lines 4–6 / text
considérons le \deleted{$\widetilde{\mathscr J}$ qu'elle \unsure{déf.}} sous-grpe jacobien \unsure{qui il|qu'il} définit, qui est donc aussi le centralisateur de $\mathfrak h$ dans $G$.
% @ lines 6–7, left edge / margin note
% slanted
\margin{à vérifier ! \uline{il faut} \uuline{car. \unsure{rés.|p.} $\ne2$}}
% @ lines 7–10 / text
\deleted{Ceci signifie} Comme $\mathscr J$ est \uline{commutatif} il centralise $\mathfrak h:\operatorname{Lie}\mathscr J$, donc $\mathscr J\subset\widetilde{\mathscr J}$, d'où égalité puisque $\widetilde J$ est connexe et $J,\widetilde J$ sont lisses de $=$ dim rel. $1$\ldots
% archive sheet number: 55
\scanpage{57}
% @ line 1, upper right / text
Dictionnaire lié à
% @ line 2, upper right / display
\[
 \mathrm{\unsure{GP|QP}}(1)\simeq\mathrm{SO}(3)
\]
% archive sheet number: 56
\scanpage{58}
% @ lines 1–7 / text
$S$ schéma. On va définir des équivalences entre les \deleted{\unsure{sept}}\inserted{neuf} catégories suivantes attachées à $S$ -- et pour $S$ variable, ce seront des équivalences mutuelles de catégories fibrées sur $\mathrm{Sch}_{/S}$ (en fait, des gerbes \unsure{fpqc|fppf|fppc} sur $\mathrm{Sch}_{/S}$), formant un ``syst. transitif d'équivalences à isom can près'' :
% @ lines 8–26, two-column dictionary and return arrow / diagram
\diagram{\begin{tikzcd}[ampersand replacement=\&,column sep=0.9em,row sep=0.5em,cells={anchor=west}]
% item line 8: torsors
 {C_{0_S}=\operatorname{Tors}(\mathrm{GP}(1)_S)} \&
 {\parbox{7cm}{catégorie des torseurs \inserted{à droite} $\unsure{P|R}$ sous $\mathrm{GP}(1)_S$}}\\
% item lines 9–12: forms of groups, three rows
 {\left\{\begin{array}{l}
 C_{1_S}=\operatorname{For}(\mathrm{GP}(1)_S)\\[1.3em]
 C'_{1_S}=\operatorname{For}(\mathrm{Sl}(2)_S)\\[1em]
 \deleted{C''_{1_S}=\operatorname{For}(\mathrm{Gl}(2)_S)}
 \end{array}\right.} \&
 {\begin{array}{l}
 \parbox{7cm}{catégorie des ``formes'' $G$ (\unsure{fpqc|fppf} ou \unsure{f.ét.|f(ét)}) du \deleted{$\mathrm{GP}(1)$} schéma en \unsure{groupe|groupes} $\mathrm{GP}(1)$}\\[.5em]
 \text{catégorie des formes }\widetilde G\text{ de }\mathrm{Sl}(2)_S\\[1em]
 \deleted{\text{catégorie des formes }\unsure{\bar G|G}\text{ de }\mathrm{Gl}(2)_S}
 \end{array}}\\
% item lines 13–16: Lie algebras, stars at left
 {\begin{array}{r} \unsure{c}(*)\\[1em](*)\\[.8em]\deleted{(*)}\end{array}
 \left\{\begin{array}{l}
 C_{2_S}=\operatorname{For}(\mathfrak{gp}(1)_S)\\[1em]
 C'_{2_S}=\operatorname{For}(\mathfrak{sl}(2)_S)\\[.8em]
 \deleted{C''_{2_S}=\operatorname{For}(\mathfrak{gl}(2)_S)}
 \end{array}\right.} \&
 {\begin{array}{l}
 \parbox{7cm}{catégorie des formes $\unsure{\mathfrak g|g}$ de l'Algèbre de Lie $\mathfrak{gp}(1)_S$ de $\mathrm{GP}(1)_S$}\\
 \parbox{7cm}{\deleted{catégorie des formes} \rule[0.3ex]{1.4em}{0.3pt}\ \rule[0.3ex]{1.4em}{0.3pt}\ \rule[0.3ex]{1.4em}{0.3pt}\ $\widetilde{\mathfrak g}$ de $\mathfrak{sl}(2)_S\simeq\underline{\operatorname{Lie}}\mathrm{Sl}(2)_S$}\\
 \deleted{\text{catégorie des formes de }\mathfrak{gl}(2)_S=\underline{\operatorname{Lie}}\mathrm{Gl}(2)_S}
 \end{array}}\\
% item lines 17–19: quaternion algebras
 {C_{3_S}=\deleted{\unsure{\operatorname{Quat}(S)}}\operatorname{Quat}(\mathcal O_S)} \&
 {\parbox{7cm}{catégorie des Algèbres d'Azumaya\inserted{$M$} de rang $4$ sur $S$ (ou Algèbres de Quaternions)}}\\
% item lines 19–23: projective lines and quadrics
 {\left\{\begin{array}{l}
 C_{5_S}=\operatorname{Drpr}(S)\\[2.6em]
 C'_{5_S}=\operatorname{Quad}(S)
 \end{array}\right.} \&
 {\begin{array}{l}
 \parbox{7cm}{catégorie des fibrés en droites projectives $X$ sur $S$}\\[1em]
 \parbox{7cm}{catégorie des couples $(P,X)$ d'un fibré en plans projectifs sur $S$ et d'une quadrique lisse $X\subset P$}
 \end{array}}\\
% item lines 24–26: special orthogonal modules
 {\left[C_{4_S}=\overset{\operatorname{Modspquad}(S)^3}{\deleted{\unsure{\operatorname{Spéc}(S)_3}}}\right.}
 \arrow[uu,bend left=44,start anchor=west,end anchor=west] \&
 {\parbox{7cm}{catégories des ``Modules spéciaux orthogonaux de rang $3$ sur $S$'' $(E,q,\omega)$, $E$ fibré vect. de rg $=3$, $\omega$ \deleted{\unsure{forme}} base de $\det(\underline E)$, $q$ forme \unsure{pol$^{\mathrm{le}}$|quad.}, $\unsure{d(q)|\delta(q)}=-\omega^{\otimes2}$}}
\end{tikzcd}}
% archive sheet number: 57
\scanpage{59}
% @ lines 1–8 / text
Ici le signe $(*)$ en marge indique que l'équivalence des 3 catégories en question avec les 8 autres n'est valable que si $S$ est premier à $2$ -- en général on a seulemt des foncteurs qui vont de chacune des huit autres catégories dans chacune des trois marquées du signe $(*)$, mais ce ne sont pas des foncteurs pl. fidèles ....

% @ lines 9–16 / text
Un dictionnaire complet devrait comprendre, notamment, la description explicite des $9\times8$\inserted{$=72$} équivalences de catégories, sans compter les \deleted{$\unsure{11.10.9=990}$}\inserted{$9.8.7=504$} isomorphismes de transitivité (si on voulait être regardant ...) On fera un peu moins, mais on essaiera de donner une image claire de la situation.

% @ lines 17–21 / text
\uline{Principe général}. On va définir des représentants typiques des onze catégories dans le cas où $S_0=\operatorname{Spec}\mathbb Z$, et on définira une opération de $G_{0\deleted{\unsure{S}}}=\mathrm{GP}(1)$ sur chacun. On établira que ceci établit une isom.
% @ line 22 / display
\[
 \mathrm{GP}(1)\simeq\operatorname{Aut}(\unsure{\mathscr X|\mathcal X|\mathfrak X}_0)
\]
% @ lines 23–24, continues / text
de $\mathrm{GP}(1)$ avec le schéma des autom. de chacun des objets $\mathscr X$ envisagés -- \unsure{en se|on se|car si}
% archive sheet number: 58
\scanpage{60}
% @ lines 1–3 / text
bornant cependant à \deleted{Spec} $S'_0=\deleted{\unsure{\mathrm{Sp}}}S-\{2\}$ dans les trois cas $(*)$ où l'opération \inserted{sur $S_0$} est fidèle, mais l'hom de grpes envisagé n'est pas épi.
% @ lines 4–12 / text
Comme les 11 catégories \inserted{fibrées} envisagées sont des gerbes (\unsure{fpqc|fppf} disons, ou ét), il s'ensuit que chacune d'elles est can. équiv. à la cat. des torseurs \inserted{à droite} sous le $\underline{\operatorname{Aut}}(\unsure{\mathscr X|\mathcal X|\mathfrak X}_0)$ correspondant, donc aussi à la cat. des torseurs à dr. sous $G_0$ (avec le grain de sel $*$). Ceci établira l'\uline{existence} d'un multidictionnaire, mais ne \inserted{nous} dispensera pas de la tâche de le développer en détail.

% @ line 13, heading / text
I Représentants typiques des onze cat. fibrées.
% @ lines 14–22, three columns / diagram
\diagram{$\begin{array}{@{}lll@{}}
% item line 14: right torsor
 1)\ C_{0_{S_0}}: & \multicolumn{2}{l}{\deleted{\unsure{R_0=|P_0=}}\text{Le torseur trivial }\inserted{R_0=(G_{S_0})_{\rm d}}\text{ sous }\mathrm{GP}(1)_{S_0}=G_{S_0}}\\[3pt]
% item lines 15–17: trivial group forms
 2)\ C_{1_{S_0}}: & \text{La forme triviale} &G_0=G_{S_0}\inserted{=\mathrm{GP}(1)}\text{ de }G_{S_0}\\
 \deleted{3)}\ C'_{1_{S_0}} &\rule[.4ex]{3.3cm}{.4pt} & \widetilde G_0=\unsure{\mathrm{Sl}|\mathrm{Gl}}(2)\text{. de }\mathrm{Sl}(2)\\
 \deleted{4)\ C^{\unsure{\prime|\prime\prime}}_{1_{S_0}}}&\deleted{\rule[.4ex]{3.3cm}{.4pt}}&\deleted{\unsure{\bar G_0}=\mathrm{Gl}(2)\text{ de }\mathrm{Gl}(2)}\\[4pt]
% item lines 18–20: trivial Lie forms
 \unsure{3|5})\ \left\{\begin{array}{l}C_{2_{S_0}}\\ C'_{2_{S_0}}\\ \deleted{C''_{\unsure{2|1}_{S_0}}}\end{array}\right.
 &\begin{array}{l}\rule[.4ex]{3.3cm}{.4pt}\\ \rule[.4ex]{3.3cm}{.4pt}\\\deleted{\rule[.4ex]{3.3cm}{.4pt}}\end{array}
 &\begin{array}{l}\unsure{\mathfrak g|g}_0=\mathfrak{gp}(1)\text{ de }\mathfrak{gp}(1)\\
 \widetilde{\mathfrak g}_0=\mathfrak{sl}(2)\text{ de }\mathfrak{sl}(2)\\
 \deleted{\unsure{\mathfrak L_0|\mathfrak g''_0}=\mathfrak{gl}(2)\text{ de }\mathfrak{gl}(2)}\end{array}\\[4pt]
% item line 21: matrix algebra
 4)\ C_{3_{S_0}}: &\multicolumn{2}{l}{\text{l'algèbre d'Azumaya triviale }M_0=\mathbb M_2(\mathcal O_{S_0})}\\
% item line 22: struck projective representative and beginning of following row
 \deleted{5)\left\{\begin{array}{l}C_{4_{S_0}}:\\ C'_{4_{S_0}}:\end{array}\right.}
 &\multicolumn{2}{l}{\deleted{X_0=\mathbb P^1_{S_0}=\mathbb P^1(\mathcal O_{S_0}^2)=\check{\mathbb P}(\mathcal O_{S_0}^2)}} 
\end{array}$}
% archive sheet number: 59
\scanpage{61}
% @ lines 1–3, row 5 / text
5) $C_{4_{S_0}}:\quad E_0=\mathcal O_{S_0}^3,\quad q_0(x,y,z)=xy+z^2,\quad\omega=e_1\wedge e_2\wedge e_3$ où $(e_1,e_{\unsure{2|3}},e_3)$ est la base canonique (on écrit $(x,y,z)$ pour $xe_1+ye_2+ze_3$)
% @ line 4, row 6 / display
\[
 6)\quad C_{5_{S_0}}:\qquad X_0=\mathbb P^1_{S_0}=\mathbb P(\mathcal O_{S_0}^2)=\check{\mathbb P}(\mathcal O_{S_0}^2)
\]
% @ lines 5–6 / text
$C'_{5_{S_0}}:\quad P_0=\mathbb P^2_{S_0}=\check{\mathbb P}(E_0),\quad C_0=\mathbb V(q_0)$ quadrique définie par la forme $q_0$ de 5°).

% @ line 7, heading / text
II \uline{Opérations de $G_0$}
% @ line 8 / text
1) Translations à gauche
% @ lines 9–11, struck block / deletion
\deleted{\unsure{3|2}) $G_0$ opère sur $\unsure{\mathfrak g|g}_0=\underline{\operatorname{Lie}}(G_0)$ par repr. adjointe. Les op. de $\widetilde G_0$ et $\unsure{\mathscr L_0|L_0|\mathcal L_0}$ sur leurs Algèbres de Lie respectives}
% @ lines 12–15 / text
2) $G_0$ opère sur lui-$=$ par autom. int$^{\mathrm s}$, les op. de $\widetilde G_0$ \deleted{\unsure{et $\mathscr L_0$}} sur \deleted{\unsure{eux|lui}}-mêmes se factorisent par les centres, $Z_0$ \deleted{\unsure{et $G_{m,S_0}$} et} on a des \deleted{isogénies} iso can
% @ line 16 / display
\[
 \widetilde G_0/Z_0=\mathrm{Sl}(2)/\mu_2\simeq\mathrm{GP}(1)=G_0\tag*{(1)}
\]
% @ line 17, crossed formula / display
\[
 \deleted{\mathscr L_0/G_{m,S_0}=\mathrm{Gl}(2)/G_m\simeq\mathrm{GP}(1)=G_0}
\]
% @ line 18 / text
d'où opérations de $G_0$ sur $\widetilde G_0$. \deleted{\unsure{et $\mathscr L_0$}}
% @ lines 19–20 / text
3) Les op. de $G_0$ sur $G_0$ et $\widetilde G_0$ définissent des op. sur $\mathfrak g_0=\underline{\operatorname{Lie}}(G_0)$ et $\widetilde{\mathfrak g}_0=\operatorname{Lie}(\widetilde G_0)$.
% @ lines 21–23 / text
4) L'opération régulière de $\check{\mathbb V}(M_0)^*$ \inserted{$\simeq\mathrm{Gl}(2)$} sur \deleted{\unsure{$M$}} \deleted{\unsure{est}}\inserted{$M_0$} est \deleted{nulle}\inserted{triviale} sur le centre $\simeq G_m\subset\mathrm{Gl}(\unsure{2|1})$, donc passe au quotient en une opération de
% @ line 24 / display
\[
 \mathrm{Gl}(2)/G_m\overset{\mathrm{dfn}}{=}\mathrm{GP}(1)\tag*{(2)}
\]
% @ line 25 / text
sur $M_0$.
% archive sheet number: 60
\scanpage{62}
% @ lines 1–5 / text
5) Considérons la forme \deleted{déterminant}\inserted{$\xi\mapsto-\det(\unsure{\vert|\xi})$} sur $M_0=\mathbb M_2(\mathcal O_{S_0})$, et considérons sa restriction au sous-Module $\mathfrak{sl}(2)_{S_0}\subset M_0$ des matrices de trace nulle, \deleted{\unsure{laquelle|le groupe}} i.e. des matrices
% @ line 6 / display
\[
 \xi=\begin{pmatrix}z&x\\y&-z\end{pmatrix}\tag*{(3)}
\]
% @ line 7 / text
donnant
% @ lines 8–9 / display
\[
 \begin{aligned}
 &\deleted{q(\xi)\overset{\mathrm{df}}{=}\det\xi=}\\
 &q_0(\xi)\overset{\mathrm{dfn}}{=}-\det\xi=xy+z^2
 \end{aligned}\tag*{(4)}
\]
% @ lines 10–14 / text
Cette forme \inserted{$-\det u$ sur \deleted{\unsure{$M_0$}} $M_0=\mathbb M_2(\mathcal O_{S_0})$} est évidemment invariante par l'action de $\mathrm{GP}(1)$, déduite par passage au quotient de celle de $\mathrm{Gl}(2)$ -- car
% @ line 15 / display
\[
 \det(g\xi g^{-1})=\det\xi\tag*{(5)}
\]
% @ lines 16–18 / text
a fortiori sa restriction à $\deleted{\unsure{\mathfrak L_0}}\widetilde{\unsure{\mathfrak g|g}}_0=\mathfrak{sl}(2)_{S_0}$ l'est. Identifiant $\widetilde{\mathfrak g}_0$ à $E_0=\mathcal O_{S_0}^3$ par
% @ lines 19–20 / display
\[
 \left\{\begin{aligned}
 (x,y,z)&\longrightarrow\begin{pmatrix}z&x\\y&-z\end{pmatrix}\\
 E_0=\mathcal O_{S_0}^3&\xrightarrow{\sim}\widetilde{\mathfrak g}_0=\mathfrak{sl}(2)_{\unsure{S|S_0}}
 \end{aligned}\right.\tag*{(6)}
\]
% @ lines 21–22 / text
on trouve une action de $G_0=\mathrm{GP}(1)$ sur $(E_0,q_0)$. Il est immédiat
% archive sheet number: 61
\scanpage{63}
% @lines 1–4 / text
aussi, comme $G_0$ opère trivialement sur \uline{tt} module inversible (car $\underline{\operatorname{Hom}}(\mathrm{GP}(1),G_m)=0$) qu'il opère trivialement sur $\underline{\det}(\widetilde{\unsure{\mathfrak g|g}}_0)\simeq\det(E_0)$, donc invarie $\omega$ en plus de $q$.

% @lines 5–9 / text
6) $\mathscr L_0=\mathrm{Gl}(2)$ opérant sur $\mathcal O_{S_0}^2$, opère sur $\mathbb P^1=\mathbb P(\mathcal O_{S_0}^{\unsure{2|?}})$, et comme $G_m$ opère trivialement, on en déduit par passage au quotient une action de $\inserted{G_0=}\mathrm{GP}(1)$ grâce à l'iso (2).

% @lines 10–11 / text
Enfin, l'opération de $G_0$ sur $(E_0,q_0)$ définit une opération sur

% @line 12 / display
\[
(P_0=\check{\mathbb P}(E_0), C_0=\mathbb V(q_0)).
\]
% @below line 12 / text
\noindent\hspace*{.24\linewidth}\tikz[baseline]{\draw (0,0) .. controls (.3,.09) and (.45,-.08) .. (.8,.01) .. controls (1.2,.09) and (1.4,-.04) .. (1.8,.01) .. controls (2.5,.02) and (2.8,-.05) .. (3.2,.04) .. controls (3.5,.08) and (3.8,.0) .. (4.1,.02);}

% @lines 13–14 / text
Nous admettrons \inserted{(pour ne pas nous éterniser...)} que tous les hom obtenus

% @left of lines 13–15 / margin note
\margin{Expliciter le thm d'isomorphie}
% slanted

% @line 15 / display
\[
G_0=\mathrm{GP}(1)\longrightarrow\underline{\operatorname{Aut}}(\unsure{\mathscr X|\mathcal X|\mathfrak X}_0)\tag*{(7)}
\]

% @lines 16–19 / text
sont des \uline{iso}, ce qui en principe fonde le dictionnaire que nous voulons. Mais nous allons \deleted{\illegible{1}} expliciter \inserted{\unsure{de ce dictionnaire|ce dictionnaire|ou dictionnaire}} les points essentiels.

% archive sheet number: 62
\scanpage{64}
% @lines 1–3 / text
III Le foncteur allant de $C_{0S}=\operatorname{Tors}(\mathrm{GP}(1)_S)$ vers les huit autres catégories est le foncteur ``torsion''

% @line 4 / display
\[
R\longmapsto\deleted{\unsure{\mathscr X|X}_0}\,R\mathbin{\wedge}_{\mathrm{GP}(1)_S}\unsure{\mathscr X|X|\mathfrak X}_{0S}\tag*{(8)}
\]

% @lines 5–10 / text
où $X_0$ est le représentant typique \inserted{sur $\mathbb Z$} de la catégorie fibrée envisagée. Le thm d'isomorphie précédent équivaut à dire que ces foncteurs sont des équiv. de cat, avec le grain de sel \deleted{\unsure{pour}} relatif aux cat. $(*)$.

% @line 11 / text
Dans le cas

% @lines 12–13 / diagram
\diagram{\begin{tikzcd}[ampersand replacement=\&,row sep=small]
C_{0S}\arrow[r]\arrow[d,equal] \& C_{5S}\arrow[d,equal]\\
\operatorname{Tors}(\mathrm{GP}(1)_S) \& \mathrm{Drpr}(S)
\end{tikzcd}}

% @lines 14–15 / text
on peut expliciter aussi ce foncteur comme

% @line 16 / display
\[
R\longmapsto R/B_{0S}\tag*{(9)}
\]

% @line 17 / text
où

% @line 18 / display
\[
B_0\subset \mathrm{GP}(1)=G_0\tag*{(10)}
\]

% @lines 19–21 / text
est le Borel-type de $\mathrm{GP}(1)$, \deleted{\unsure{groupe} stabilisant le point} \inserted{déduit par passage au quotient par $\mu_2$ du} \deleted{matrices} groupe des matrices

% @line 22 / display
\[
\widetilde B_0=\left\{\begin{pmatrix}\lambda&\mu\\ 0&\lambda^{-1}\end{pmatrix}\ \middle|\ \begin{aligned}&\lambda\deleted{\unsure{,\mu}}\text{ dans }G_m\\&\mu\text{ dans }G_a\end{aligned}\right\}\tag*{(10)}
\]

% @lines 23–24 / text
\deleted{\unsure{on en voit|on en conclut}} qui est aussi le stabilisateur de la droite $\mathcal O_S e_1$ dans $\deleted{\unsure{E}}\mathcal O_S^2$ -- s'identifie

% archive sheet number: 63
